\appendix

\section{Additional Research Resources}
\label{appendix:additional-resources}

This appendix provides additional resources related to the implementation and
reproducibility of the research presented in this thesis. The resources include
the source code of the proposed models and the software implementation used
for model deployment. These materials are provided to facilitate further
investigation, reproduction, and extension of the proposed approach by
researchers and practitioners.

\subsection{Source Code}

The complete source code developed for this research is publicly available
through the following GitHub repository:

\begin{center}
\href{https://github.com/chhaytheanly/ultralytics}
{\textbf{GitHub Repository: Ultralytics -- Research Implementation}}
\end{center}

The repository contains the source code, configuration files, and other
implementation resources associated with the experiments presented in this
thesis. Researchers interested in reproducing or extending the proposed
method can clone the repository and inspect the corresponding implementation.

In particular, the custom neural network modules developed for this research
are located in the following directory:

\begin{verbatim}
ultralytics/nn/modules/
\end{verbatim}

This directory contains the implementations of the customized neural network
components used in the proposed architecture. The corresponding classes and
functions define the structure and behavior of the modules incorporated into
the model. Examining these implementations provides additional technical
details that complement the architectural description presented in the main
chapters of this thesis.

\subsection{Model Engine Package}

In addition to the research implementation, the developed model engine has
been packaged and published as a Python package to facilitate its deployment
and reuse. The package is publicly available through the Python Package Index
(PyPI):

\begin{center}
\href{https://pypi.org/project/pyfarm-detection/}
{\textbf{PyPI Package: pyfarm-detection}}
\end{center}

The package provides the software components required to deploy and utilize
the developed detection model in practical applications. It is currently
under active development, with ongoing improvements intended to enhance its
functionality, usability, and performance.

The source code of the model engine is also available through the following
GitHub repository:

\begin{center}
\href{https://github.com/chhaytheanly/pyfarm.git}
{\textbf{GitHub Repository: pyfarm}}
\end{center}

Given the ongoing development of the model engine, users, researchers, and developers are encouraged to refer to the repository for the latest updates, documentation,
and usage instructions. One important note is that I am looking for collaborators to help improve the model engine and its associated resources. 
If you are interested in contributing to the development of the model engine, please feel free to pull requests, report issues, or reach out to me through the repository's contact channels.
Your contributions and feedback are highly appreciated and will help enhance the overall quality and usability of the model engine.\\


These resources are provided to support reproducibility and encourage further
development of the proposed system. Researchers and developers may use the
provided implementations as a reference for further experimentation,
evaluation, and extension of the presented work.
